\documentclass[11pt]{article}
\title{On the Closure of Compositional Hierarchies in Digital Symmetries}
\author{A rigorous researcher-engineer bridging abstract mathematics and concrete digital systems.}
\usepackage[margin=1in]{geometry}
\usepackage{graphicx}
\usepackage{amsmath,amssymb,mathtools}
\usepackage{hyperref}
\graphicspath{{media/}{media/diagrams/}{images/}{artwork/}}

\begin{document}

\section*{Symmetry Groups of Boolean Functions and Circuits}

This chapter studies symmetry as an organizing principle for Boolean functions and circuits. The central objects are actions of the symmetric group \(\Sigma_n\) on \(n\)-variable Boolean functions, their automorphism groups, and the induced orbit--stabilizer structure that partitions functions into symmetry classes. These notions provide a bridge between abstract equivalence relations on truth tables and concrete transformations of circuit descriptions.

Let \(f : \mathbb{B}^n \to \mathbb{B}\) be a Boolean function. A permutation \(\pi \in \Sigma_n\) acts on inputs by reindexing coordinates, and we write \(f^\pi\) for the transported function obtained by permuting variables according to \(\pi\). The automorphism group of \(f\) is the subgroup
\[
\operatorname{Aut}(f) = \{\pi \in \Sigma_n : f^\pi = f\},
\]
which measures the internal symmetries of the function. The orbit of \(f\) under \(\Sigma_n\) consists of all functions obtainable by permuting variables, while the stabilizer records the permutations that leave the function unchanged. In circuit settings, the same ideas apply to the input interface of a circuit, and symmetry can be used to simplify search, canonicalization, and equivalence checking.

A finer notion of equivalence is NPN equivalence, which combines negation of inputs, permutation of inputs, and negation of the output. Two Boolean functions are NPN-equivalent if one can be transformed into the other by these operations. This equivalence relation is especially useful in classification problems, because it identifies functions that differ only by trivial relabelings and polarity changes. Canonical forms provide a representative for each NPN class, enabling systematic classification and comparison. In practice, a canonicalizer attempts to map each function to a unique normal form, while preserving the equivalence class under the chosen transformation group.

For circuit analysis, NPN equivalence is often paired with structural invariants extracted from the truth table or from spectral data. A canonical form can serve as a compact fingerprint for a function or subcircuit, supporting database indexing, duplicate detection, and benchmark generation. The chapter's artifact-oriented perspective treats canonicalization not only as a theoretical classification tool but also as an executable component in a symmetry analysis pipeline.

Beyond permutation symmetries, Boolean functions admit linear and affine symmetries. The general linear group \(\mathrm{GL}(n,2)\) acts on \(\mathbb{B}^n\) by invertible linear transformations over \(\mathbb{F}_2\), and the affine general linear group \(\mathrm{AGL}(n,2)\) extends this action by translations. These groups are relevant when one studies symmetries that preserve algebraic structure rather than merely variable names. Linear and affine equivalence are coarser than permutation equivalence and are especially important in contexts such as cryptography, coding theory, and reversible or quantum-flavored computation.

Spectral methods provide a complementary approach to symmetry analysis. Using the Walsh--Hadamard transform, a Boolean function can be represented by its spectral coefficients, which encode correlations with parity functions. These coefficients often reveal hidden regularities that are not obvious from the truth table alone. In particular, approximate symmetry can be quantified by a spectral correlation measure \(\varepsilon(r)\), where \(r\) indexes a candidate symmetry or transformation and the value of \(\varepsilon(r)\) indicates how closely the function aligns with that symmetry. Such measures are useful when exact invariance fails but near-symmetry still carries algorithmic or structural significance.

The spectral viewpoint is especially effective for detecting approximate automorphisms, ranking candidate permutations, and guiding search over large symmetry groups. In a workflow for symmetry discovery, one may first compute spectral signatures, then restrict attention to transformations that preserve or nearly preserve those signatures, and finally verify exact equivalence or automorphism conditions by direct evaluation. This combination of coarse spectral filtering and exact checking is often more efficient than exhaustive enumeration.

The chapter also emphasizes the relationship between symmetry and classification. Functions and circuits can be organized into equivalence classes under \(\Sigma_n\)-actions, NPN transformations, or affine transformations, depending on the level of abstraction required. Each choice of group action induces a different notion of canonical representative and a different notion of complexity for the classification problem. The resulting hierarchy of symmetry notions mirrors the broader theme of the book: closure properties and compositional structure become more tractable when the relevant invariances are made explicit.

An important practical outcome is the construction of a symmetry oracle and canonicalizer with benchmarks. A symmetry oracle is a procedure that, given a Boolean function or circuit, returns information about its automorphism group, candidate symmetries, or equivalence class membership. A canonicalizer uses this information to produce a normal form. Benchmarks then evaluate correctness, speed, and robustness across families of functions with varying degrees of symmetry. Such artifacts make the chapter's ideas reproducible and testable, and they provide a concrete bridge from theory to implementation.

In summary, the chapter develops three intertwined perspectives on symmetry in Boolean settings: group actions and automorphisms, NPN and affine equivalence for classification, and spectral methods for exact and approximate symmetry detection. Together these tools support both mathematical understanding and practical algorithms for analyzing Boolean functions and circuits.


\end{document}
